\documentclass[12pt]{article}

\pagestyle{empty}
\setcounter{secnumdepth}{2}
%----------------------------------------------------------------------------------------
%   Packages and configurations
%----------------------------------------------------------------------------------------
\usepackage[dutch]{babel}
\usepackage{hyperref} %Allow for references
\hypersetup{
colorlinks,
citecolor=black,
filecolor=black,
linkcolor=black,
urlcolor=black
} %Set up hyperlink colours.
\newcommand{\sectionbreak}{\clearpage} % Should start every section on its own page
\usepackage{geometry} % Required to change the page size to A4
\geometry{a4paper} % Set the page size to be A4 as opposed to the default US Letter
\usepackage{chngpage}
\usepackage{appendix}
%REMOVE IF HEADER/FOOTER BROKEN
\usepackage{fancyhdr} % Required for custom headers
\usepackage{extramarks} % Required for headers and footers
\usepackage{lastpage} % Required to determine the last page for the footer
%-----------------1
\topmargin=0cm
\oddsidemargin=0cm
\textheight=22.0cm
\textwidth=16cm
\parindent=0cm
\parskip=0.15cm
\topskip=0truecm
\raggedbottom
\abovedisplayskip=3mm
\belowdisplayskip=3mm
\abovedisplayshortskip=0mm
\belowdisplayshortskip=2mm
\normalbaselineskip=12pt
\normalbaselines
\usepackage{listings}
\usepackage[svgnames,table,xcdraw]{xcolor}
\lstset { 
language=C,
frame=single,
escapeinside={\%*}{*)}, 
breaklines=true,  
backgroundcolor=\color{black!5},
basicstyle=\footnotesize,
commentstyle=\color{mygreen},
numberstyle=\tiny\color{mygray},
rulecolor=\color{black},
keywordstyle=\color{blue},
}
\definecolor{mygreen}{rgb}{0,0.6,0}
\definecolor{mygray}{rgb}{0.5,0.5,0.5}
\definecolor{mymauve}{rgb}{0.58,0,0.82}
\usepackage{wasysym}
\usepackage{lipsum}
\pagestyle{fancy}
\lhead{Guus \textsc{Beckett} \& Jim \textsc{van Abkoude}} % Top left header
\rhead{Eindrapport} % Top center header
\lfoot{\includegraphics[height=1cm]{avans}} % Bottom left footer
\cfoot{} % Bottom center footer
\rfoot{Pagina\ \thepage} % Bottom right footer
\renewcommand\headrulewidth{0.4pt} % Size of the header rule
\renewcommand\footrulewidth{0.4pt} % Size of the footer rule
\usepackage{graphicx} % Required for including pictures

\usepackage{float} % Allows putting an [H] in \begin{figure} to specify the exact location of the figure
\usepackage{wrapfig} % Allows in-line images such as the example fish picture
\usepackage{lipsum} % Used for inserting dummy 'Lorem ipsum' text into the template
\usepackage{pdfpages}
\usepackage[font={footnotesize}]{caption}
\graphicspath{ {../Images/Logos/}{../Images/} }
\setlength\parindent{0pt} % Removes all indentation from paragraphs
%\usepackage{showframe}
\newcommand*{\SignatureAndDate}[1]{%
    \par\noindent\makebox[2.5in]{\hrulefill} \hfill\makebox[2.0in]{\hrulefill}%
    \par\noindent\makebox[2.5in][l]{#1}      \hfill\makebox[2.0in][l]{Date}%
}%Signature package
\begin{document}
\begin{titlepage}
\pagenumbering{Roman}
\newcommand{\HRule}{\rule{\linewidth}{0.5mm}} % Defines a new command for the horizontal lines, change thickness here

\center % Center everything on the page

\includegraphics[height=3cm] {avans}\\% Include a department/university logo - this will require the graphicx package
\textsc{\Large Avans Hogeschool Breda}\\[0.5cm] % Major heading such as course name
\textsc{\large Netwerktechnologie}\\[0.5cm] % Minor heading such as course title
\HRule \\[0.4cm]
{ \huge \bfseries Eindrapport}\\[0.4cm] % Title of your document
\HRule \\[1.5cm]

\begin{minipage}{0.4\textwidth}
\begin{flushleft} \large
\emph{Auteurs:}\\
Guus \textsc{Beckett} \\% Your name 
Jim \textsc{van Abkoude}
\end{flushleft}
\end{minipage}
~
\begin{minipage}{0.4\textwidth}
\begin{flushright} \large
\emph{Docenten:} \\
Andries \textsc{van Dongen} \\ % Supervisor's Name
Johan \textsc{Talboom}
\end{flushright}
\end{minipage}\\[4cm]

{\large \today}\\[3cm] % Date, change the \today to a set date if you want to be precise

Versie: 0.1.0

\vfill % Fill the rest of the page with whitespace

\end{titlepage}

\clearpage
\section*{Voorwoord}
\addcontentsline{toc}{section}{Voorwoord}
Dit eindrapport zal beschrijven bla bla bla, Zeroshell bleek kut om te PXE booten, bla bla bla.
\\
Guus Beckett \& Jim van Abkoude \\
\today \\
Breda
\newpage
\tableofcontents
\newpage
\pagenumbering{arabic}
\section{Requirements}
De requirements voor dit project waren verdeeld over de 3 apparaten die in het netwerk moeten komen en bla bla
\subsection{Gateway} % (fold)
\label{sub:gateway}
Een van de door school gestelde eisen was het plaatsten van een gateway tussen het Avans hogeschool netwerk en de TSLR van dit project, dit was waarschijnlijk bedoeld als bescherming voor het Avans netwerk tegen mogelijke fouten die de student zou kunnen maken en het omzeilen van MAC adress authenticatie. Origineel was er in dit project sprake van nog een vierde apparaat in het project, namelijk de imageserver. De imageservers van Avans waren niet beschikbaar gesteld aan de studenten en dus werden de verantwoordelijkheden van de imageserver naar de gateway verplaatst.
\\De gateway moest in staat zijn als zowel een DHCP en een DNS server te fungeren waarin de DNS server een eigen TLD moest hosten voor binnenin het lokale netwerk. Ook moest gateway packetfiltering toepassen.
\\Aangezien de verantwoordelijkheden van de imageserver zijn verwerkt in de gateway vallen de requirements van de imageserver nu onder de gateway. De gateway moest daarom ook een TFTP-server, NFS-server en PXE-bootserver hebben. Het imageserver deel had ook enkele bonus requirements, dit waren het draaien van een webserver met PHP en het draaien van een MySQL server op de gateway.
% subsection gateway (end)
\subsection{Thin Server Linux Router} % (fold)
\label{sub:thin_server_linux_router}
Het hoofddeel van het project was het installeren van een TSLR (Thin Server Linux Router) deze via PXE te laten booten waarbij hij zijn bootfiles verkrijgt via TFTP, een remote filesystem te laten gebruiken via NFS en een gebruik te maken van een static WAN IP.
\\De TSLR moest naast PXE ook firewall en routingfuncties hebben. De routerfuncties van de TSLR gaan over het beheren van een LAN, een DMZ en het routeren van packets tussen deze 2 netwerken.
\\Zeer veel firewall verantwoordelijkheden moeten ook door de TSLR uitgevoerd worden zoals het blokeren van nieuwe connecties vanuit de DMZ naar het LAN, het blokkeren van nieuwe connecties van WAN naar LAN, het droppen van alle ICMP packets, het toelaten van alle verbindingen van LAN naar WAN. het toelaten van alle verbindingen van LAN naar DMZ en het instellen van static routing.
\\Er waren ook een paar overige requirements zoals het systeem logs remote opslaan en een SSH shell laten draaien. Naast de overige requirements waren er ook bonus requirements, dit waren het draaien van een webserver en een infopagina met status en NIC statistieken.
% subsection thin_server_linux_router (end)
\section{Vooronderzoek}
Tijdens het vooronderzoek lag de focus op het vinden van de juiste OS voor de juiste apparaten. Dit vooronderzoek bestond uit 3 delen. De eisen duidelijk opstellen, de OS bepalen van elk van de drie apparaten en een planning opstellen.\\

De eisen die naar voren zijn gekomen tijdens dit onderzoek zijn terug te vinden in hoofdstuk~\ref{sub:thin_server_linux_router} Thin Server Linux Router op pagina~\pageref{sub:thin_server_linux_router}.\\

Het kiezen van de juiste OS voor het juiste apparaat werd gedaan door middel van het gebruik van virtuele machines, waarin de verschillende OS's die geschikt leken uitvoerig werden getest. Aangezien er tijdens het vooronderzoek geen beschikking was over de hardware waarop de OS's uiteindelijk moesten gaan draaien, waren de keuzes die tijdens deze fase van het onderzoek zijn gedaan niet gericht op de hardware waarop de OSsen moesten draaien, maar alleen op de features van de gesteste OS's. De OS's die gekozen waren na dit vooronderzoek zijn:
\begin{description}
\item[Gateway - OpenBSD] De beveiligingsgerichte focus van OpenBSD maakt het uitermate geschikt om de gateway te zijn tussen het Internet en het lokale netwerk.
\item[TSLR - Zeroshell] Het feit dat Zeroshell gemaakt is om te fungeren als routersoftware maakte het een voordehandliggende keus om te gebruiken op de TSLR.
\item[DMZ server - SuperbMiniServer] De OS die op de DMZ server zou draaien moest gemakkelijk op te zetten zijn en in te stellen zijn, vandaar de keus voor deze distro, die focust op het gemakkelijk maken van het managen van een server.
\end{description}
Ten slotte is de planning gemaakt. De planning die opgesteld is focuste vooral op de volgorde van het installeren van de OS's op de verschillende apparaten. Allereerst zou de TSLR zijn OS krijgen, waarna de Gateway ge\"installeerd worden. Als laatste zou de DMZ server zijn OS krijgen. Hierna zou de PXE boot opgezet worden en de oplevering zijn van het project.
\section{Ontwerp}
\section{Implementatie}
\section{Testen}
\end{document}
